% Paso 1: Plantilla de Observación Inicial - VERSIÓN CORREGIDA
% Proyecto: Análisis Exegético Profundo
% Ministerio "La Gloria es del Señor"

\documentclass[11pt,a4paper]{article}
\usepackage[utf8]{inputenc}
\usepackage[spanish]{babel}
\usepackage[margin=2.5cm]{geometry}
\usepackage{graphicx}
\usepackage{fancyhdr}
\usepackage{xcolor}
\usepackage{tcolorbox}
\usepackage{enumitem}
\usepackage{tabularx}
\usepackage{multirow}
\usepackage{amsfonts}
\usepackage{fontspec}
\usepackage{titlesec}

% Configuración de colores (basados en el diseño web)
\definecolor{forestgreen}{RGB}{27,67,50}
\definecolor{sagegreen}{RGB}{45,90,39}
\definecolor{olivegreen}{RGB}{64,145,108}
\definecolor{mintgreen}{RGB}{82,183,136}
\definecolor{lightmint}{RGB}{116,198,157}
\definecolor{softcream}{RGB}{216,243,220}
\definecolor{papergray}{RGB}{118,128,150}

% Configuración de encabezados y pies de página
\pagestyle{fancy}
\fancyhf{}
\fancyhead[L]{\textcolor{forestgreen}{\textbf{Análisis Exegético Profundo}}}
\fancyhead[R]{\textcolor{papergray}{Paso 1: Observación Inicial}}
\fancyfoot[C]{\textcolor{papergray}{\thepage}}
\fancyfoot[R]{\textcolor{papergray}{\scriptsize Ministerio "La Gloria es del Señor"}}

% Configuración de títulos
\titleformat{\section}
{\color{forestgreen}\Large\bfseries}
{\thesection}{1em}{}

\titleformat{\subsection}
{\color{sagegreen}\large\bfseries}
{\thesubsection}{1em}{}

% Configuración de cajas de texto
\tcbuselibrary{skins}
\newtcolorbox{infobox}[1]{
    colback=softcream,
    colframe=olivegreen,
    fonttitle=\bfseries,
    title=#1,
    left=10pt,
    right=10pt,
    top=10pt,
    bottom=10pt
}

\newtcolorbox{checklistbox}{
    colback=white,
    colframe=mintgreen,
    left=15pt,
    right=15pt,
    top=10pt,
    bottom=10pt,
    boxrule=2pt
}

\begin{document}

% Título principal
\begin{center}
    {\Huge\color{forestgreen}\textbf{PASO 1}}\\[0.5cm]
    {\Large\color{sagegreen}\textbf{OBSERVACIÓN INICIAL}}\\[0.3cm]
    {\large\color{papergray}Plantilla para Análisis Exegético}\\[1cm]
    
    \begin{tcolorbox}[colback=olivegreen!10, colframe=olivegreen, width=0.8\textwidth]
        \centering
        \textcolor{forestgreen}{\textbf{Proyecto: Análisis Exegético Profundo}}\\
        \textcolor{papergray}{Metodología de 8 Pasos para Estudio Bíblico Riguroso}
    \end{tcolorbox}
\end{center}

\vspace{1cm}

% Información del estudiante
\begin{infobox}{Información del Proyecto}
\begin{tabularx}{\textwidth}{|X|X|}
\hline
\textbf{Nombre del Estudiante:} & \underline{\hspace{6cm}} \\[0.8cm]
\hline
\textbf{Texto Bíblico Seleccionado:} & \underline{\hspace{6cm}} \\[0.8cm]
\hline
\textbf{Fecha de Inicio:} & \underline{\hspace{6cm}} \\[0.8cm]
\hline
\textbf{Congregación/Iglesia:} & \underline{\hspace{6cm}} \\[0.8cm]
\hline
\end{tabularx}
\end{infobox}

\section{Objetivo del Paso 1: Observación Inicial}

La observación inicial es el fundamento de todo análisis exegético riguroso. En este primer paso, nos acercamos al texto bíblico con mente abierta y ojos atentos, buscando ver lo que realmente dice antes de interpretarlo.

\subsection{Descripción del Proceso}

Lee el texto múltiples veces en diferentes traducciones. Identifica personajes, lugares, eventos y palabras que se repiten. Anota tus primeras impresiones sin hacer interpretaciones profundas todavía. El objetivo es \textbf{observar} antes de \textbf{interpretar}.

\begin{infobox}{Principio Hermenéutico Fundamental}
\textit{"La observación cuidadosa precede a la interpretación correcta. Lo que el texto dice debe establecerse antes de determinar lo que significa."} — Metodología Exegética Clásica
\end{infobox}

\section{Lista de Verificación}

\begin{checklistbox}
\textbf{Marca cada elemento completado:}

\begin{itemize}[leftmargin=1cm]
    \item[$\square$] Leer el texto en al menos 3 traducciones diferentes
    \item[$\square$] Identificar y marcar palabras y frases clave
    \item[$\square$] Marcar conectores lógicos (por tanto, pero, porque, sin embargo, etc.)
    \item[$\square$] Anotar preguntas iniciales que surgen del texto
    \item[$\square$] Observar la estructura del párrafo y divisiones lógicas
    \item[$\square$] Identificar personajes, lugares y eventos mencionados
    \item[$\square$] Notar repeticiones, contrastes y paralelos
    \item[$\square$] Registrar primeras impresiones sin interpretar
\end{itemize}
\end{checklistbox}

\newpage

\section{Hoja de Trabajo: Observación Textual}

\subsection{Traducciones Consultadas}

\begin{tabularx}{\textwidth}{|c|X|X|}
\hline
\textbf{\#} & \textbf{Traducción} & \textbf{Observaciones Específicas} \\
\hline
1 & \underline{\hspace{4cm}} & \\[1.5cm]
\hline
2 & \underline{\hspace{4cm}} & \\[1.5cm]
\hline
3 & \underline{\hspace{4cm}} & \\[1.5cm]
\hline
4 & \underline{\hspace{4cm}} & \\[1.5cm]
\hline
\end{tabularx}

\subsection{Palabras y Frases Clave Identificadas}

\begin{tabularx}{\textwidth}{|X|X|X|}
\hline
\textbf{Palabra/Frase} & \textbf{Versículo} & \textbf{Observación Inicial} \\
\hline
& & \\[1cm]
\hline
& & \\[1cm]
\hline
& & \\[1cm]
\hline
& & \\[1cm]
\hline
& & \\[1cm]
\hline
& & \\[1cm]
\hline
\end{tabularx}

\subsection{Conectores Lógicos y Estructura}

\textbf{Conectores encontrados en el texto:}
\vspace{2cm}

\textbf{Estructura observada del pasaje:}
\vspace{3cm}

\subsection{Personajes, Lugares y Eventos}

\begin{tabularx}{\textwidth}{|X|X|X|}
\hline
\textbf{Personajes} & \textbf{Lugares} & \textbf{Eventos/Acciones} \\
\hline
& & \\[2cm]
\hline
\end{tabularx}

\newpage

\subsection{Repeticiones, Contrastes y Paralelos}

\textbf{Palabras o conceptos que se repiten:}
\vspace{2cm}

\textbf{Contrastes observados:}
\vspace{2cm}

\textbf{Paralelos o comparaciones:}
\vspace{2cm}

\subsection{Preguntas Iniciales}

Lista las preguntas que surgen naturalmente de tu observación del texto:

\begin{enumerate}[leftmargin=1cm]
    \item \underline{\hspace{12cm}}
    \vspace{0.5cm}
    \item \underline{\hspace{12cm}}
    \vspace{0.5cm}
    \item \underline{\hspace{12cm}}
    \vspace{0.5cm}
    \item \underline{\hspace{12cm}}
    \vspace{0.5cm}
    \item \underline{\hspace{12cm}}
    \vspace{0.5cm}
    \item \underline{\hspace{12cm}}
    \vspace{0.5cm}
    \item \underline{\hspace{12cm}}
    \vspace{0.5cm}
    \item \underline{\hspace{12cm}}
    \vspace{0.5cm}
\end{enumerate}

\subsection{Primeras Impresiones}

\textbf{¿Cuál parece ser el tema principal del pasaje?}
\vspace{2cm}

\textbf{¿Qué tono o ambiente transmite el texto?}
\vspace{2cm}

\textbf{¿Qué aspectos del texto te llaman más la atención?}
\vspace{2cm}

\begin{infobox}{Recordatorio Importante}
En esta fase, \textbf{NO} busques interpretar o explicar el significado profundo. Simplemente observa y registra lo que ves en el texto. La interpretación vendrá en los pasos posteriores.
\end{infobox}

\section{Siguiente Paso}

Una vez completada esta observación inicial, estarás listo para proceder al \textbf{Paso 2: Análisis Léxico}, donde estudiaremos el significado de las palabras clave en los idiomas originales.

\vspace{1cm}

\begin{center}
\begin{tcolorbox}[colback=mintgreen!20, colframe=mintgreen, width=0.7\textwidth]
\centering
\textcolor{forestgreen}{\textbf{¡Felicitaciones!}}\\
\textcolor{papergray}{Has completado el primer paso hacia un análisis exegético profundo.}
\end{tcolorbox}
\end{center}

\end{document}